% Tapis de Quillen, with Quillen's offprint and preprint (cote 162-5, typed part): archive cote 162-5, site page 58 (PDF page 59).
% preprint. Draft compared: none.
% AI-transcribed by Claude Opus 5.5 (claude-opus-5-5, adaptive thinking, effort high), Message Batches API, batch t162_full (msgbatch_01UAh3zyWPd454AeDkhnmQ33), 2026-11-13; not yet checked by eye.
\scanpage{162-5/58}
$f:E - 0 \rightarrow F - 0$ which is of degree a power of $p$ on each fiber.
When $E$ is a line bundle, then $\hand{\psi}^pE = E^{\otimes p}$ and $f(v) = v^{\otimes p}$ is
the desired map to prove the conjecture for $E$ . Unfortunately
when $E$ has dimension greater than one, then $\hand{\psi}^pE$ is only a vir-%
tual bundle so the method fails. However if $E$ is an algebraic
vector bundle (locally free sheaf of finite rank) over a prescheme
$V$ of characteristic $p$ , then in virtue of the Frobenius homomor-%
phism there is an algebraic vector bundle $E^{(p)}$ whose class in
$K_{alg}(V)$ represents $\hand{\psi}^pE$ .\hand{*} If $E$ is given by a cocycle $(g_{ij})$ of
$n\mathrm{x}n$ matrices with respect to an open covering of $V$ , then $E^{(p)}$
is given by the cocycle $(g_{ij}^{(p)})$ , where $g_{ij}^{(p)}$ denotes the matrix
whose entries are the p-th powers of the corresponding entries of
$g_{ij}$ . Moreover there is a canonical map $F:E \rightarrow E^{(p)}$ given in lo-%
cal coordinates by $(X_i) \rightarrow (X_i^p)$ which is of (inseparable) degree
$p^{\hand{\dim E}}$ on each fiber. Hence the hypotheses for Adams' criterion is
verified in the algebraic situation and it remains to reduce the
topological situation to the algebraic one.

It evidently suffices to prove the conjecture when $X$ is the
Grassmannian of complex $N$ planes in $n + N$ space and where
$x \in K(X)$ is the class of the canonical quotient bundle. Let $G_Z$
be the Grassmannian scheme over Spec Z representing the functor
``isomorphism classes of n-dimensional vector bundles with $n + N$
generating sections'' and let $Q_Z$ be the canonical n-dimensional
bundle. Then $X = G_{C,cl}$ and $x$ is the class of $Q_{C,cl}$ , where

\hand{* That $E^{(p)}$ represents $\psi^pE$ follows since $E \mapsto E^{(p)}$ defines a
homomorphism $K_{alg} \rightarrow K_{alg}$ which coincides with $\psi^p$ on line bundles,
together with the splitting principle: $K_{alg}(PE) \unsure{\cong|=} K_{alg}(X)[T]/(\lambda_{-T}(E))$.}

\mspage{58}
